\sectiontheme{rustaccent}
\section{Additional Language Targets}
\label{sec:additional-targets}

Every remaining client target takes the same \texttt{.emi.yml} through
its own sub-compiler, with the same \texttt{--path}/\texttt{--output}
pair every other Emi compiler already uses. Each one below is
client-only --- no server, no CLI, no entity/gorm persistence --- and
most share the same two housekeeping tags: \texttt{no-sdk} (skip
embedding the runtime files) and \texttt{no-pkg} (skip the
language-native package manifest --- \texttt{requirements.txt},
\texttt{pom.xml}, \texttt{.csproj}, \texttt{composer.json},
\texttt{pubspec.yaml}, \texttt{README.md}, depending on the target).

\subsection{Python}
\texttt{dataclasses}-based DTOs and an \texttt{httpx} client, synchronous
by default:
\begin{lstlisting}
emi python --path module.emi.yml --output ./out/python
emi python --path module.emi.yml --output ./out/python --tags async
\end{lstlisting}
\texttt{--tags async} swaps every generated client function to
\texttt{async def}, backed by \texttt{httpx.AsyncClient}, instead.

\subsection{Dart}
DTO classes with hand-generated \texttt{toJson}/\texttt{fromJson} and a
\texttt{package:http} client:
\begin{lstlisting}
emi dart --path module.emi.yml --output ./out/dart
\end{lstlisting}

\subsection{C\#}
\texttt{System.Text.Json}-based DTOs and an \texttt{HttpClient} transport:
\begin{lstlisting}
emi csharp --path module.emi.yml --output ./out/csharp
\end{lstlisting}

\subsection{Java}
Jackson-based DTOs on top of \texttt{java.net.http.HttpClient}:
\begin{lstlisting}
emi java --path module.emi.yml --output ./out/java
\end{lstlisting}

\subsection{PHP}
Reflection-based DTO hydration and a \texttt{curl} HTTP transport:
\begin{lstlisting}
emi php --path module.emi.yml --output ./out/php
\end{lstlisting}

\subsection{C}
Vendored \texttt{cJSON} (de)serialization and \texttt{libcurl} transport
--- no nullable value types without pointers, a documented scope limit:
\begin{lstlisting}
emi c --path module.emi.yml --output ./out/c
\end{lstlisting}

\paragraph{ESP-IDF server mode.} \texttt{--tags esp-server} flips the C target from
client to \emph{server}: every action header then carries a typed handler
the developer implements, plus \texttt{<action>\_register(server, ...)} on top of
\texttt{esp\_http\_server} (routing, path/query/body decoding, JSON and status
handling are generated). \texttt{method: reactive} actions become websocket
sessions with \texttt{on\_open}/\texttt{on\_message}/\texttt{on\_close} callbacks and
typed \texttt{\_send()}/\texttt{\_broadcast()}. \texttt{--tags mongoose-server} generates the
same action headers on top of Mongoose's event loop instead, which runs on
macOS, Linux, Windows and ESP-IDF (bring your own \texttt{mongoose.c}/\texttt{.h}).
See \texttt{examples/esp-server}.
\begin{lstlisting}
emi c --path module.emi.yml --output ./main/emi --tags esp-server
\end{lstlisting}

\paragraph{SQLite entities.} \texttt{--tags sqlite} turns the module's entities into C structs with
storage on \texttt{sqlite3.h}: \texttt{Entity\_insert/\_get/\_update/\_patch/\_delete/\_list/\_query} and a
\texttt{<module>\_schema\_apply(db)} that creates and extends the tables. The tables and the
row layout come from the same code as \texttt{emi entity-sql --tags sqlite}, so the C code and the
SQL script agree on every column. Scalars, flattened objects and arrays of objects are stored.
\texttt{Entity\_query} takes an \texttt{EmiSqlQuery}: a json-logic filter (the format of the Go Browse
action; values are bound, unknown fields are refused), a sort, \texttt{start\_index} and
\texttt{items\_per\_page}, and returns the matching entities with the total. \texttt{Entity\_patch}
changes only the fields of a json object.

\paragraph{Packed entities, without SQLite.} \texttt{--tags pack} gives every entity
\texttt{Entity\_pack} and \texttt{Entity\_unpack} (and \texttt{\_unpack\_prefix} for several in a row): a
compact binary string for a file, an NVS blob or a message, where there is no SQLite. It works alone
(no \texttt{sqlite3.h}) or together with \texttt{sqlite}, and uses the same layout. No field names or
types are stored: a null/bool/``needs 8 bytes'' flag pair per field, integers as zigzag varints, floats
as 4 bytes when that is exact, strings with a length. A 190-byte json entity packs to about 55 bytes.
Fields may only be appended to an entity which has been stored: a newer program reads an older blob
(missing fields keep their defaults), an older program refuses a newer one, and damaged or truncated input
is refused without reading out of bounds. A header compiled with \texttt{-DEMI\_NO\_SQLITE} leaves the
SQLite functions out. A handler of the server targets can say why it fails with
\texttt{emihttpd\_set\_error("...")} before returning \texttt{NULL} and a status of 400 or more.

\subsection{C++: generic and Unreal Engine}
One compiler, two dialects, picked with \texttt{--dialect} (or the
equivalent \texttt{--tags unreal}, handy when a caller is already
building a \texttt{--tags} list): portable ISO C++17 for desktop/ESP-IDF/
Arduino by default, or \texttt{USTRUCT}/\texttt{UPROPERTY} DTOs usable
from Blueprint when the Unreal dialect is selected.
\begin{lstlisting}
emi cpp --path module.emi.yml --output ./out/cpp                  # generic ISO C++17
emi cpp --path module.emi.yml --output ./out/cpp-unreal \
  --dialect unreal                                                # Unreal Engine
\end{lstlisting}

\begin{pullquote}
\texttt{--path}, \texttt{--output}, and \texttt{--tags} is the interface
every Emi compiler shares, this section's seven targets included --- run
\texttt{emi tags} at any time to print the full, current tag list for
every registered target, straight from the compiler's own source of
truth rather than a copy that can drift out of date here.
\end{pullquote}
